\documentclass[10pt,a4paper]{amsart}
\usepackage[margin=19mm]{geometry}
\usepackage{amsmath,amssymb,mathtools}
\usepackage[scr=rsfs,cal=euler]{mathalfa}
\usepackage{xfrac}
\usepackage[unicode,hidelinks,bookmarks=false]{hyperref}
\newcommand{\ie}{i.e.\ }
\newcommand{\cf}{cf.\ }
\newcommand{\resp}{respectively\ }
\newcommand{\Subsection}{\subsection}
\providecommand{\nobreakdash}{\nobreak\hbox{-}}
\setlength{\emergencystretch}{2em}
\multlinegap0pt
\usepackage{bm}
\usepackage{bbm}
\usepackage{dsfont}
\usepackage{xspace}
\usepackage{cancel}
\usepackage{mathtools}
\usepackage{amsthm}
\makeatletter
\@ifundefined{lemma}{\newtheorem{lemma}{Lemma}}{}
\makeatother
\newcommand{\E}{\mathbb{E}}
\newcommand{\bA}{\boldsymbol{A}}
\newcommand{\bC}{\boldsymbol{C}}
\newcommand{\bD}{\boldsymbol{D}}
\newcommand{\bH}{\boldsymbol{H}}
\newcommand{\bI}{\boldsymbol{I}}
\newcommand{\bM}{\boldsymbol{M}}
\newcommand{\bP}{\boldsymbol{P}}
\newcommand{\bR}{\boldsymbol{R}}
\newcommand{\bS}{\boldsymbol{S}}
\newcommand{\bh}{\boldsymbol{h}}
\newcommand{\bigo}{\mathcal{O}}
\newcommand{\bzero}{\boldsymbol{0}}
\newcommand{\bzet}{\boldsymbol{\zeta}}
\newcommand{\calH}{\mathcal{H}}
\newcommand{\pnlim}{\xrightarrow[n \to \infty]{P}}
\title[Robust Spatial Confounding Adjustment via Basis Voting]{Robust Spatial Confounding Adjustment via Basis Voting}
\author{Anik Burman, Elizabeth L. Ogburn, Abhirup Datta}
\date{Source: arXiv:2510.22464v2 (CC BY 4.0)}
\begin{document}
\maketitle

\noindent\emph{Source and licence.} Robust Spatial Confounding Adjustment via Basis Voting. Version arXiv:2510.22464v2 (CC BY 4.0), distributed under the Creative Commons Attribution 4.0 International licence, as recorded in the arXiv licence metadata for this exact version and shown on the abstract page https://arxiv.org/abs/2510.22464v2. Full rights notice: licenses/arxiv-2510.22464-CC-BY-4.0.txt.\par\smallskip

\noindent\textit{Editorial scope.} Counted unit: the Lemma whose source label is \texttt{lem:A-lim} in the article (source0.tex lines 1200--1206, source label \texttt{lem:A-lim}), delivered together with the article's own complete original proof of it (source0.tex lines 1207--1322, one \texttt{proof} environment of 116 lines). No mathematics is written, repaired or re-derived: statement and proof are the article's own text line for line. Required context retained uncounted: lem:dot-prod-lim (lines 831--837); prop:tight-matrix-inv (lines 861--867); eq:proj-simplify (lines 1190--1198). Every definition, display and label of those lines is retained unchanged. Omitted from this package and not redistributed: 1--830, 838--860, 868--1189, 1199--1199, 1323--1504. Nothing inside a retained excerpt is omitted or altered: every retained body is delivered exactly as the article's lines are written, so no omission receipt of any kind accompanies this package. Citations: the retained excerpts carry no citation command at all, so no bibliography entry of the article is transcribed and no reference is redistributed. Reference adaptation: none was needed and none was made; every reference inside the delivered unit resolves inside it, so no unresolved reference or citation is produced. Printed slips kept literal: no printed word, symbol, hypothesis or number of the counted statement or of its proof was altered. Layout and class: the article's own class is \texttt{article} with packages including amsmath, amssymb, hyperref, iftex, comment, enumerate, fontenc, inputenc, textcomp, lmodern, parskip, xcolor, longtable, booktabs; the wrapper is the lane's pinned \texttt{amsart} editorial wrapper at 10pt on A4, which restates the theorem-like environments the retained text needs and adds the article's own macro definitions that the retained text needs (\texttt{\textbackslash E}, \texttt{\textbackslash bA}, \texttt{\textbackslash bC}, \texttt{\textbackslash bD}, \texttt{\textbackslash bH}, \texttt{\textbackslash bI}, \texttt{\textbackslash bM}, \texttt{\textbackslash bP}, \texttt{\textbackslash bR}, \texttt{\textbackslash bS}, \texttt{\textbackslash bh}, \texttt{\textbackslash bigo}, \texttt{\textbackslash bzero}, \texttt{\textbackslash bzet}, \texttt{\textbackslash calH}, \texttt{\textbackslash pnlim}), with extra wrapper packages (bm, bbm, dsfont, xspace, cancel, mathtools). The delivered numbering is the unit's own. Assets: no retained span contains a figure environment, a graphics-inclusion command, a PDF-page inclusion, a table or a TikZ picture; no image file is copied into this package, the package declares no asset dependency and no image file is redistributed with it. Licence: version arXiv:2510.22464v2 of this article is distributed under the Creative Commons Attribution 4.0 International licence, as recorded in the arXiv licence metadata for this exact version and shown on the abstract page https://arxiv.org/abs/2510.22464v2. A full rights notice is delivered at licenses/arxiv-2510.22464-CC-BY-4.0.txt. Characters: every retained excerpt is pure ASCII, so the delivered engine, XeLaTeX with its default Latin Modern text font, reports no missing character.
\begin{lemma}\label{lem:dot-prod-lim}
    For any basis function $h_j$ from the orthogonal class $\calH$, defining its evaluation at coordinates $S_1, S_2,\cdots,S_n \overset{iid}{\sim} f_S>0$ as $\bh_{n,j} = (h_j(S_1), \cdots,h_j(S_n))^\top$ and for iid noise vector $\bzet_n = (\zeta_1, \cdots, \zeta_n) \sim N(\bzero,\sigma^2_\zeta\bI_n)$ where $\bzet_n \perp \bS_n = (S_1, \cdots, S_n)^\top$, the following holds:
    \begin{enumerate}
        \item\label{basis-basis-dot} $\dfrac{1}{n} \bh_{n,j}^\top\bh_{n,l} \pnlim \mathbb{I}(j = l)$ where $\mathbb{I}(\cdot)$ is the indicator function for any $j,l \geq 1$
        \item\label{basis-noise-dot} $\dfrac{1}{n} \bh_{n,j}^\top\bzet_n \pnlim 0$ for any $j \geq 1$
    \end{enumerate}
\end{lemma}

\begin{lemma}\label{prop:tight-matrix-inv}
    Let $\bA_n$ be a $k \times k$ finite dimensional matrix of the form $\bA_n = \frac{1}{n}\sum_{l=1}^n A_\ell$ such that each of $A_\ell$ are iid square matrices of dimension $k$ such that $\E A_\ell = \bA$ with $||\bA||_2<\infty$. Then we have:
$$
\left|\left|\sqrt{n}\left(\bA_n^{-1} - \bA^{-1}\right)\right|\right|_2 = \bigo_P(1)
$$
where $||\cdot||_2$ is the spectral norm.
\end{lemma}

\begin{equation}\label{eq:proj-simplify}
    \begin{split}
        & \frac{1}{n}\bP_d(\bI_n - \bP_{-j,d})\bP_d\\
        = & \frac{1}{n}\bH_{d}\left(\bH_{d}^\top\bH_{d}\right)^{-1}\bH_{d}^\top(\bI_n - \bP_{-j,d})\bH_{d}\left(\bH_{d}^\top\bH_{d}\right)^{-1}\bH_{d}^\top\\
        = & \frac{1}{n}\bH_{d}\left(\bH_{d}^\top\bH_{d}\right)^{-1}\bH_{d}^\top\left(\bI_n - \bH_{-j,d}\left(\bH_{-j,d}^\top\bH_{-j,d}\right)^{-1}\bH_{-j,d}^\top\right)\bH_{d}\left(\bH_{d}^\top\bH_{d}\right)^{-1}\bH_{d}^\top\\
        = & \frac{1}{n}\bH_{d}\underbrace{\left(\frac{1}{n}\bH_{d}^\top\bH_{d}\right)^{-1}\left(\frac{1}{n}\bH_{d}^\top\bH_{d} - \frac{1}{n}\bH_{d}^\top\bH_{-j,d}\left(\frac{1}{n}\bH_{-j,d}^\top\bH_{-j,d}\right)^{-1}\frac{1}{n}\bH_{-j,d}^\top\bH_{d}\right)\left(\frac{1}{n}\bH_{d}^\top\bH_{d}\right)^{-1}}_{\bA_{j,d}}\frac{1}{n}\bH_{d}^\top\\
        = &  \frac{1}{n}\bH_{d} \bA_{j,d}\frac{1}{n}\bH_{d}^\top
    \end{split}
\end{equation}

\begin{lemma}\label{lem:A-lim}
    For the $d \times d$ matrix $\bA_{j,d}$ defined in Equation \ref{eq:proj-simplify}, we have:
    \begin{itemize}
        \item $\bA_{j,d} \pnlim \bA^{lim}_{j,d}\coloneqq diag(\left\{\mathbb{I}(u = j), 1 \leq u \leq d\right\})$
        \item $\left|\left|\sqrt{n}\left(\bA_{j,d} - \bA^{lim}_{j,d}\right)\right|\right|_2 = \bigo_P(1)$ where $||\cdot||_2$ denotes the spectral norm.
    \end{itemize}
\end{lemma}

\begin{proof}
    Using Lemma \ref{lem:dot-prod-lim}, we can say that $\frac{1}{n}\bH_{d}^\top\bH_{d} \pnlim \bI_d$ and $\frac{1}{n}\bH_{-j,d}^\top\bH_{-j,d} \pnlim \bI_{d-1}$, and similarly using the same lemma, we can say that:
    \begin{equation}\label{eq:cross-basis-inner}
        \begin{split}
              \frac{1}{n}\bH_{-j,d}^\top\bH_{d} = & \frac{1}{n}\begin{bmatrix}
            \bh^\top_{n,1}\\
            \vdots\\
            \bh^\top_{n,j-1}\\
            \bh^\top_{n,j+1}\\
            \vdots\\
            \bh^\top_{n,d}
        \end{bmatrix}
        \begin{bmatrix}
    \bh_{n,1} & \cdots & \bh_{n,j-1}&\bh_{n,j}&\bh_{n,j+1}& \cdots& \bh_{n,d}
\end{bmatrix}\\
= & \frac{1}{n}\begin{bmatrix}
    \bH^\top_{1:(j-1),n}\\
     \bH^\top_{(j+1):d,n}\\
\end{bmatrix}_{(d-1)\times n}\begin{bmatrix}
    \bH_{1:(j-1),n} & \bh_{n,j} & \bH_{(j+1):d,n}
\end{bmatrix}_{n \times d}\\
\pnlim & \begin{bmatrix}
    \bI_{j-1} & \bzero_{j-1} & \bzero_{(j-1) \times (d-j)}\\
    \bzero_{(d-j)\times (j-1)} & \bzero_{d-j} & \bI_{d-j} 
\end{bmatrix} = \bR_{j,d}
        \end{split}
    \end{equation}
with:
\begin{equation}\label{eq:drop-inner-drop}
    \bR_{j,d}^\top\bR_{j,d} = \begin{bmatrix}
    \bI_{j-1} & \bzero_{j-1} & \bzero_{(j-1) \times (d-j)}\\
    \bzero^\top_{j-1} & 0 & \bzero_{d-j}\\
    \bzero_{(d-j) \times (j-1)} & \bzero_{d-j} & \bI_{d-j}\\
\end{bmatrix}
\end{equation}
Hence, we have:
$$
\begin{aligned}
    \bA_{j,d} = & \left(\frac{1}{n}\bH_{d}^\top\bH_{d}\right)^{-1}\left(\frac{1}{n}\bH_{d}^\top\bH_{d} - \frac{1}{n}\bH_{d}^\top\bH_{-j,d}\left(\frac{1}{n}\bH_{-j,d}^\top\bH_{-j,d}\right)^{-1}\frac{1}{n}\bH_{-j,d}^\top\bH_{d}\right)\left(\frac{1}{n}\bH_{d}^\top\bH_{d}\right)^{-1}\\
    \pnlim & \bI^{-1}_d\left(\bI_d - R^\top_{j,d}\bI_{d-1}^{-1}R_{j,d}\right)\bI^{-1}_d\\
    = & \bI_d - \begin{bmatrix}
    \bI_{j-1} & \bzero_{j-1} & \bzero_{(j-1) \times (d-j)}\\
    \bzero^\top_{j-1} & 0 & \bzero_{d-j}\\
    \bzero_{(d-j) \times (j-1)} & \bzero_{d-j} & \bI_{d-j} \\
\end{bmatrix} = diag(\left\{\mathbb{I}(u = j), 1 \leq u \leq d\right\}) = \bA^{lim}_{j,d}
\end{aligned}
$$
This completes the proof of the first part lemma. Next, we move to proving the uniform tightness of a matrix deviation. 

We will be using Lemma \ref{prop:tight-matrix-inv} in proving the second part of our lemma. Recall how We defined the matrix $\bA_{j,d}$ as:
$$
A_{j,d} = \left(\frac{1}{n} \bH_d^\top \bH_d \right)^{-1} \left( \frac{1}{n} \bH_d^\top \bH_d - \frac{1}{n} \bH_d^\top \bH_{-j,d} \left( \frac{1}{n} \bH_{-j,d}^\top \bH_{-j,d} \right)^{-1} \frac{1}{n} \bH_{-j,d}^\top \bH_d \right) \left(\frac{1}{n} \bH_d^\top \bH_d \right)^{-1}.
$$
and its probability limit $\bA_{j,d}^{lim} = \text{diag}(\mathbb{I}(u = j), 1 \leq u \leq d)$ which is a $d \times d$ matrix with $1$ at position $(j,j)$ and $0$ elsewhere. Let us define:
$$
\begin{aligned}
    \bM_n \coloneqq & \frac{1}{n} \bH_d^\top \bH_d \pnlim \bI_d\\
\bC_n^\top \coloneqq & \frac{1}{n} \bH_d^\top \bH_{-j,d} \pnlim \bR^\top_{j,d} \text{ (Eqn. \ref{eq:cross-basis-inner})}\\
\bD_n \coloneqq &\frac{1}{n} \bH_{-j,d}^\top \bH_{-j,d} \pnlim \bI_{d-1}    
\end{aligned}
$$
Then,
$$
\bA_{j,d} = \bM_n^{-1} \left(\bM_n - \bC_n^\top \bD_n^{-1} \bC_n\right) \bM_n^{-1} = \bM_n^{-1} - \bM_n^{-1}\bC_n^\top \bD_n^{-1} \bC_n \bM_n^{-1}
$$
Also, observe that we can write down $\bA_{j,d}^{lim} = \bI_d - \bR^\top_{j,d}\bR_{j,d}$ where $\bR^\top_{j,d}\bR_{j,d}$ has all the diagonal entries to be 1 except the $j$-th diagonal element which is 0 as seen in Equation \ref{eq:drop-inner-drop}. We denote $\bR^*_{j,d}\coloneqq \bR^\top_{j,d}\bR_{j,d}$. Thus, we have:
\begin{equation}
    \begin{split}
        & \sqrt{n}\left(\bA_{j,d} - \bA^{lim}_{j,d}\right)\\
        = & \sqrt{n}\left[\left(\bM_n^{-1} - \bI_d\right) - \left(\bM_n^{-1}\bC_n^\top \bD_n^{-1} \bC_n \bM_n^{-1} - \bR^*_{j,d}\right)\right]\\
        = & \sqrt{n}\left(\bM_n^{-1} - \bI_d\right) - \sqrt{n}\left[\bM_n^{-1}(\bC_n^\top \bD_n^{-1} \bC_n - \bR^*_{j,d} + \bR^*_{j,d})\bM_n^{-1} - \bR^*_{j,d}\right]\\
        = & \!\begin{aligned}[t]
        & \sqrt{n}\left(\bM_n^{-1} - \bI_d\right)\\
        - & \bM_n^{-1}\sqrt{n}\left(\bC_n^\top \bD_n^{-1} \bC_n - \bR^*_{j,d}\right)\bM_n^{-1}\\
        - &\sqrt{n}\left(\bM_n^{-1}\bR^*_{j,d}\bM_n^{-1} - \bR^*_{j,d}\right) \\
    \end{aligned}\\
    = & \!\begin{aligned}[t]
        & \sqrt{n}\left(\bM_n^{-1} - \bI_d\right)\\
        - & \bM_n^{-1}\left[\bC_n^\top \sqrt{n}(\bD_n^{-1} - \bI_{d-1}) \bC_n + \sqrt{n}(\bC_n^\top\bC_n - \bR^*_{j,d})\right]\bM_n^{-1}\\
        - &\sqrt{n}\left[(\bM_n^{-1}-\bI_d + \bI_d)\bR^*_{j,d}(\bM_n^{-1}-\bI_d + \bI_d) - \bR^*_{j,d}\right] \\
    \end{aligned}\\
    = & \!\begin{aligned}[t]
        & \sqrt{n}\left(\bM_n^{-1} - \bI_d\right)\\
        - & \bM_n^{-1}\left[\bC_n^\top \sqrt{n}(\bD_n^{-1} - \bI_{d-1}) \bC_n + \sqrt{n}(\bC_n^\top\bC_n - \bR^*_{j,d})\right]\bM_n^{-1}\\
        - &\left[\sqrt{n}(\bM_n^{-1}-\bI_d)\bR^*_{j,d}(\bM_n^{-1}-\bI_d)+\sqrt{n}(\bM_n^{-1}-\bI_d)\bR^*_{j,d} + \bR^*_{j,d}\sqrt{n}(\bM_n^{-1}-\bI_d)\right] \\
    \end{aligned}
    \end{split}    
\end{equation}
Unless otherwise specified, we will refer to $||\cdot||$ as the spectral norm for matrices. Thus, using these simplifications, we have:
\begin{equation}
    \begin{split}
        & \left|\left|\sqrt{n}\left(\bA_{j,d} - \bA^{lim}_{j,d}\right)\right|\right|_2\\
        \leq & \!\begin{aligned}[t]
        & \left|\left|\sqrt{n}\left(\bM_n^{-1} - \bI^{-1}_d\right)\right|\right|_2\\
        + & \left|\left|\bM_n^{-1}\left[\bC_n^\top \sqrt{n}(\bD_n^{-1} - \bI^{-1}_{d-1}) \bC_n + \sqrt{n}(\bC_n^\top\bC_n - \bR^*_{j,d})\right]\bM_n^{-1}\right|\right|_2\\
        + &
        \left|\left|\left[\sqrt{n}(\bM_n^{-1}-\bI^{-1}_d)\bR^*_{j,d}(\bM_n^{-1}-\bI^{-1}_d)+\sqrt{n}(\bM_n^{-1}-\bI^{-1}_d)\bR^*_{j,d} + \bR^*_{j,d}\sqrt{n}(\bM_n^{-1}-\bI^{-1}_d)\right] \right|\right|_2\\
        \end{aligned}\\
        \leq & \!\begin{aligned}[t]
        & \left|\left|\sqrt{n}\left(\bM_n^{-1} - \bI^{-1}_d\right)\right|\right|_2\\
        + & \left|\left|\bM_n^{-1}\right|\right|^2_2\left(\left|\left|\bC_n^\top \sqrt{n}(\bD_n^{-1} - \bI^{-1}_{d-1}) \bC_n\right|\right|_2 + \left|\left|\sqrt{n}(\bC_n^\top\bC_n - \bR^*_{j,d})\right|\right|_2\right)\\
        + &
        \left|\left|\sqrt{n}(\bM_n^{-1}-\bI^{-1}_d)\bR^*_{j,d}(\bM_n^{-1}-\bI^{-1}_d)\right|\right|_2+\left|\left|\sqrt{n}(\bM_n^{-1}-\bI^{-1}_d)\bR^*_{j,d}\right|\right|_2\\ 
        + &\left|\left|\bR^*_{j,d}\sqrt{n}(\bM_n^{-1}-\bI^{-1}_d)\right|\right|_2\\
        \end{aligned}\\
              \leq & \!\begin{aligned}[t]
        & \left|\left|\sqrt{n}\left(\bM_n^{-1} - \bI^{-1}_d\right)\right|\right|_2\\
        + & \left|\left|\bM_n^{-1}\right|\right|^2_2\left(\left|\left|\bC_n \right|\right|^2_2\left|\left|\sqrt{n}(\bD_n^{-1} - \bI^{-1}_{d-1})\right|\right|_2 + \left|\left|\sqrt{n}(\bC_n^\top\bC_n - \bR^*_{j,d})\right|\right|_2\right)\\
        + &
        \left(\left|\left|\sqrt{n}(\bM_n^{-1}-\bI^{-1}_d)\right|\right|_2\left|\left|\bR^*_{j,d}\right|\right|_2\left|\left|(\bM_n^{-1}-\bI^{-1}_d)\right|\right|_2+\left|\left|\sqrt{n}(\bM_n^{-1}-\bI^{-1}_d)\right|\right|_2\left|\left|\bR^*_{j,d}\right|\right|_2\right)\\
        + & \left|\left|\bR^*_{j,d}\right|\right|_2\left|\left|\sqrt{n}(\bM_n^{-1}-\bI^{-1}_d)\right|\right|_2\\
        \end{aligned}\\
    \end{split}
\end{equation}
Let us talk about the spectral norms of some of the components of the upper bound of the quantity $\left|\left|\sqrt{n}\left(\bA_{j,d} - \bA^{lim}_{j,d}\right)\right|\right|_2$. As $\bM_n$ converges to $\bI_d$ and matrix inversion and spectral norm are continuous maps on the space of matrices, hence we can say $\left|\left|\bM_n^{-1}\right|\right|^2_2 = \bigo_P(1)$. Similarly, we can say the same thing about $\left|\left|\bC_n \right|\right|^2_2$ being tight. Also note that since $\sqrt{n}(\bC_n - \bR_{j,d})$ has a matrix CLT due to its structure as $\bC_n$ can be written as sum of $n$ iid matrices of dimension $d \times (d-1)$, combined with the fact that $\bA \mapsto \bA^\top\bA$ is a continuous map, we can say $\left|\left|\sqrt{n}(\bC_n^\top\bC_n - \bR^\top_{j,d}\bR_{j,d})\right|\right|_2$ is $\bigo_P(1)$. Also from Equation \ref{eq:drop-inner-drop}, it is also clear that $\left|\left|\bR^*_{j,d}\right|\right|_2 = 1$. The remaining terms of the upper bound can be said to be $\bigo_P(1)$ by using Proposition \ref{prop:tight-matrix-inv}. This completes the proof of the second part of the lemma.
\end{proof}

\end{document}
